\subsection{Примеры символьной генерации многочленов}

В таблице \ref{tab:examples} приведены результаты работы программной реализации для различных параметров. Для полей характеристики $2$ многочлены содержат значительно меньше одночленов, а их коэффициенты равны единице, что упрощает аппаратную реализацию.

\begin{table}[h]
\centering
\caption{Параметры и размеры открытого ключа}
\label{tab:examples}
\begin{tabular}{|c|c|c|c|c|c|c|}
\hline
Поле & $\eta$ & $\mu$ & Степень & Многочленов & Одночленов & Размер ключа \\
\hline
$\mathrm{GF}(13)$ & 2 & 2 & 3 & 4 & 40 & 20 байт \\
$\mathrm{GF}(2^4)$ & 2 & 2 & 3 & 4 & 8 & 4 байта \\
$\mathrm{GF}(2^3)$ & 3 & 3 & 3 & 9 & 66 & 25 байт \\
\hline
\end{tabular}
\end{table}

\subsubsection{Пример для $\mathrm{GF}(2^4)$, $\eta=2$, $\mu=2$}

Для указанных параметров получены следующие многочлены открытого ключа (переменные $x_0, x_1, x_2, x_3$):

\begin{align*}
y_0 &= x_0^9 + x_0^3 x_2^6,\\
y_1 &= x_0^6 x_2^3 + x_2^9,\\
y_2 &= x_0^6 x_1^3 + x_0^2 x_1 x_2^4 x_3^2,\\
y_3 &= x_0^4 x_1^2 x_2^2 x_3 + x_2^6 x_3^3.
\end{align*}

\subsubsection{Пример для $\mathrm{GF}(2^3)$, $\eta=3$, $\mu=3$}

Для размерности $n=9$ получены следующие многочлены (приведены первые три):

\begin{align*}
y_0 &= x_0^9 + x_3^9 + x_6^9,\\
y_1 &= x_0^6 x_3^3 + x_0^3 x_6^6 + x_3^6 x_6^3,\\
y_2 &= x_0^6 x_6^3 + x_0^3 x_3^6 + x_3^3 x_6^6.
\end{align*}

\subsubsection{Пример для $\mathrm{GF}(13)$, $\eta=2$, $\mu=2$}

Для сравнения приведём многочлены для поля нечётной характеристики:

\begin{align*}
y_0 &= x_0^9 + 18x_0^7x_1^2 + 108x_0^5x_1^4 + 216x_0^3x_1^6 + 9x_0^3x_2^6 + 108x_0^3x_2^4x_3^2 \\
    &\quad + 324x_0^3x_2^2x_3^4 + 54x_0x_1^2x_2^6 + 648x_0x_1^2x_2^4x_3^2 + 1944x_0x_1^2x_2^2x_3^4,\\
y_1 &= 3x_0^6x_2^3 + 18x_0^6x_2x_3^2 + 36x_0^4x_1^2x_2^3 + 216x_0^4x_1^2x_2x_3^2 + 108x_0^2x_1^4x_2^3 \\
    &\quad + 648x_0^2x_1^4x_2x_3^2 + 3x_2^9 + 54x_2^7x_3^2 + 324x_2^5x_3^4 + 648x_2^3x_3^6,\\
y_2 &= 27x_0^6x_1^3 + 54x_0^4x_1^5 + 36x_0^2x_1^7 + 243x_0^2x_1x_2^4x_3^2 + 324x_0^2x_1x_2^2x_3^4 \\
    &\quad + 108x_0^2x_1x_3^6 + 8x_1^9 + 162x_1^3x_2^4x_3^2 + 216x_1^3x_2^2x_3^4 + 72x_1^3x_3^6,\\
y_3 &= 81x_0^4x_1^2x_2^2x_3 + 54x_0^4x_1^2x_3^3 + 108x_0^2x_1^4x_2^2x_3 + 72x_0^2x_1^4x_3^3 \\
    &\quad + 36x_1^6x_2^2x_3 + 24x_1^6x_3^3 + 81x_2^6x_3^3 + 162x_2^4x_3^5 + 108x_2^2x_3^7 + 24x_3^9.
\end{align*}

Как видно, многочлены над полями нечётной характеристики содержат большое количество одночленов и большие коэффициенты, что делает их менее удобными для практической реализации по сравнению с полями характеристики $2$.

\subsection{Время генерации и производительность}

Для демонстрационных параметров ($\eta=3$, $\mu=3$, $\mathrm{GF}(2^3)$, степень 3) время генерации открытого ключа составило менее $0.1$ секунды, шифрование одного вектора — около $0.001$ секунды, расшифрование — $0.005$ секунды. Для реальных параметров ($\eta=5$, $\mu=17$, $\mathrm{GF}(2^8)$, степень 257) размер открытого ключа, согласно оценкам, составляет примерно $7$ КБ, что в десятки раз меньше, чем у известных постквантовых схем (например, GeMSS требует $352$–$3040$ КБ).
